\documentclass[../BTL.tex]{subfiles}
\begin{document}

\section*{Tổng quan chương}
Chương 3 đã trình bày chi tiết về kiến trúc phần mềm, hiện thực hóa các lớp xử lý và các kết quả thực nghiệm thực tiễn của đề tài. Chương 4 này sẽ đưa ra những đánh giá tổng kết cuối cùng về quá trình thực hiện bài tập lớn môn học Thực tập cơ sở. Nội dung chương bao gồm hai phần chính: (i) Phần Kết luận tổng hợp các kết quả đạt được, ưu điểm và mặt hạn chế còn tồn tại của hệ thống; và (ii) Phần Hướng phát triển vạch ra các định hướng nghiên cứu, nâng cấp giải pháp kỹ thuật trong tương lai.

\section{Kết luận}
\label{section:4.1}
Đề tài nghiên cứu và xây dựng hệ thống tương tác người - máy bằng cử chỉ tay thời gian thực dựa trên xử lý ảnh và học máy, kết hợp mô phỏng môi trường 3D đã hoàn thành các mục tiêu chính trong phạm vi bài tập lớn. Sau quá trình thiết kế, triển khai thực tế và thử nghiệm hiệu năng, một số kết quả chính đã đạt được bao gồm:
\begin{enumerate}
    \item \textbf{Xây dựng hệ thống tương tác thời gian thực phù hợp:} Hệ thống đạt tốc độ xử lý khoảng 30 - 32 FPS trong điều kiện thử nghiệm, cho phép phản hồi tương đối tốt đối với chuyển động và cử chỉ của người dùng. Tổng độ trễ toàn trình ước lượng trong điều kiện thử nghiệm cục bộ dao động từ 35 - 45 ms, tạo cảm giác tương tác bám đuổi tương đối mượt.
    \item \textbf{Phương pháp nhận dạng cử chỉ chính xác và linh hoạt:} Kết hợp hiệu quả giữa thuật toán luật hình học Heuristics (dựa trên quy tắc co/duỗi ngón tay) và phương pháp học máy có huấn luyện (sử dụng mạng nơ-ron nhân tạo dạng truyền thẳng MLP và triển khai dưới định dạng TensorFlow Lite). Mô hình phân loại cử chỉ tĩnh đạt độ chính xác khoảng $99\%$ trên tập kiểm thử, cho thấy khả năng nhận dạng ổn định đối với các trạng thái bàn tay cơ bản. Mô hình phân loại cử chỉ động đạt độ chính xác khoảng $84\%$ trên tập kiểm thử, phản ánh khả năng nhận dạng bước đầu đối với các quỹ đạo chuyển động, tuy nhiên vẫn còn nhầm lẫn ở một số lớp có hình dạng tương đồng.
    \item \textbf{Đánh giá đối chứng thực nghiệm rõ ràng:} Kết quả so sánh trực tiếp chỉ ra mô hình học máy MLP vượt trội rõ rệt so với thuật toán Heuristics về độ chính xác toàn cục ($99.0\%$ so với $78.27\%$), F1-score và khả năng mở rộng cử chỉ mới, trong khi Heuristics chỉ giữ lợi thế về việc không tốn thời gian huấn luyện và cấu hình tối giản.
    \item \textbf{Hiện thực hóa các ứng dụng tương tác 3D trực quan:} Thiết lập cơ chế truyền dữ liệu cục bộ qua giao thức UDP để kết nối module Python AI với môi trường Unity. Xây dựng hệ xương bàn tay ảo mô phỏng \texttt{3DHandModel} tích hợp cơ chế cầm nắm, ném vật thể ảo dựa trên mô phỏng vật lý và trò chơi tương tác giải trí \texttt{FruitNinja} điều khiển bằng ngón trỏ.
    \item \textbf{Giảm hiện tượng rung giật và nâng cao độ ổn định tương tác:} Áp dụng thành công thuật toán bộ lọc thông thấp thích nghi One Euro Filter trong Fruit Ninja kết hợp với cơ chế giới hạn vùng hoạt động (\textit{Active Region Remapping}) giúp giảm thiểu tối đa hiện tượng rung lắc tọa độ của đầu ngón tay khi tương tác. Đồng thời, việc cài đặt giải thuật dung hợp cử chỉ (\textit{BothMatch}) và trễ chống nhiễu cầm nắm (\textit{Debounce}) đã giảm đáng kể tình trạng nhảy nhãn nhầm lẫn tức thời của mô hình AI, giúp trải nghiệm tương tác vật lý (cầm/nắm/ném) trong ứng dụng mô phỏng xương bàn tay đạt độ ổn định và chính xác cao trong thực tế.
\end{enumerate}

Bên cạnh những kết quả tích cực đã đạt được, hệ thống vẫn tồn tại một số hạn chế kỹ thuật cần được khắc phục: (i) độ chính xác trích xuất khớp của MediaPipe phụ thuộc nhiều vào điều kiện ánh sáng xung quanh, dễ bị nhảy nhiễu hoặc mất dấu vết (lost tracking) khi ánh sáng yếu hoặc ngược sáng; (ii) hiện tượng tự che khuất (self-occlusion) khi bàn tay xoay nghiêng hoặc các ngón tay gập khuất sau lòng bàn tay làm MediaPipe dự đoán sai tọa độ khớp, dẫn đến việc hiển thị tay ảo trong Unity bị méo lệch; và (iii) phân loại cử chỉ tĩnh đôi khi bị nhầm lẫn giữa các lớp có cấu trúc gần giống nhau khi người dùng di chuyển tay quá nhanh giữa các trạng thái.

\section{Hướng phát triển}
\label{section:4.2}
Để nâng cao độ ổn định, độ chính xác và mở rộng khả năng ứng dụng thực tiễn của hệ thống, hướng nghiên cứu phát triển tiếp theo của đề tài sẽ tập trung vào các công việc sau:
\begin{enumerate}
    \item \textbf{Nâng cấp mô hình phân loại cử chỉ động:} Do mô hình cử chỉ động hiện tại sử dụng MLP trên vector quỹ đạo đã làm phẳng nên khả năng khai thác quan hệ thời gian còn hạn chế. Trong tương lai, có thể nghiên cứu các mô hình chuỗi thời gian như LSTM, GRU hoặc Transformer để học tốt hơn đặc trưng động học của cử chỉ.
    \item \textbf{Hỗ trợ tương tác hai bàn tay và đa cử chỉ phối hợp:} Mở rộng hệ thống để trích xuất và nhận diện đồng thời cử chỉ của cả hai tay. Thiết kế các cử chỉ phối hợp nâng cao như zoom vật thể bằng hai tay, xoay đối tượng 3D hoặc tương tác giao diện người dùng (UI) hai tay để nâng cao số bậc tự do nhập liệu.
    \item \textbf{Mở rộng hệ thống tương tác hai bàn tay ở góc nhìn thứ nhất (First-Person):} Trích xuất và xử lý đồng thời 21 điểm landmark của cả hai bàn tay từ Python truyền qua giao thức UDP để điều khiển trực tiếp cặp tay ảo trong môi trường game 3D. Nghiên cứu và áp dụng cơ chế vật lý ragdoll, các ràng buộc khớp (joint limits) và vùng va chạm ảo (colliders) độc lập cho từng bàn tay. Hệ thống cho phép người chơi sử dụng cả hai tay phối hợp thực hiện các thao tác phức tạp như cầm, nắm, ném vật thể, điều khiển vũ khí (súng, kiếm) hoặc giải các câu đố đòi hỏi sự chuyển động linh hoạt của từng ngón tay. Giải pháp này hướng tới việc tái hiện trải nghiệm tương tác không gian tương tự các ứng dụng VR/AR chuyên dụng nhưng chạy trên nền tảng webcam RGB thông thường, giúp tăng khả năng tiếp cận cho người dùng.
    \item \textbf{Ứng dụng luồng dữ liệu cử chỉ và tọa độ ngón tay vào không gian game 2D:} Ánh xạ luồng dữ liệu 21 điểm landmark từ Python vào hệ tọa độ camera Orthographic hoặc Canvas UI trong không gian 2D. Hệ thống tập trung bắt các cử chỉ tĩnh/động kết hợp trích xuất vị trí độc lập của từng đầu ngón tay để tạo ra các tính năng gameplay độc đáo như: vẽ quỹ đạo để đại pháp thuật thi triển (cast spell), thay đổi trạng thái co/duỗi các ngón tay để tạo hình và xoay các khối gạch giải đố (giống game Block Blast), hoặc tận dụng khoảng cách và vị trí linh hoạt giữa các ngón tay để làm tính năng ngắm bắn, bóp cò, tương tác vi mô với môi trường 2D.
\end{enumerate}

\section*{Kết chương}
Chương 4 đã tổng kết các kết quả đạt được và ý nghĩa thực tiễn của đề tài bài tập lớn Thực tập cơ sở, đồng thời phân tích các ưu điểm cùng những hạn chế kỹ thuật hiện tại của hệ thống tương tác cử chỉ tay thời gian thực. Các hướng phát triển được đề xuất là cơ sở quan trọng để tiếp tục nâng cấp hệ thống trong tương lai, hướng tới một giải pháp tương tác người - máy tự nhiên, mượt mà và tiện dụng hơn.

\end{document}